\chapter{Eye Care}
\index{Eye Care}

\section*{Definition and Overview}
Eye care is the range of habits, examinations, and treatments that keep the eyes healthy and vision clear throughout life. The eyes allow us to learn, work, and connect with the world, and protecting them is important at every age. Good eye care includes regular eye examinations, protection from sun and injury, managing conditions that affect the eyes, and seeking prompt care for symptoms. Many causes of vision loss are preventable or treatable when detected early, which is why regular checks matter.

\section*{Epidemiology}
Vision problems are common in many countries. Refractive errors affect a large proportion of the population, age-related macular degeneration and glaucoma are leading causes of vision loss in older people, and cataract is widespread with age. Diabetic eye disease affects many people with diabetes, and a significant proportion of vision loss is avoidable with screening and treatment. Many people have not had an eye examination in recent years.

\section*{Aetiology and Pathophysiology}
Eye health is affected by many factors.
\begin{itemize}
    \item \textbf{Ageing:} cataract, macular degeneration, and glaucoma become more common with age
    \item \textbf{Refractive error:} the shape of the eye causes blurring that glasses and contacts correct
    \item \textbf{Chronic disease:} diabetes damages the retina, and high blood pressure affects the eye's blood vessels
    \item \textbf{Sun exposure:} ultraviolet light increases the risk of cataract and eye surface cancers
    \item \textbf{Injury:} eye injuries occur at work, in sport, and around the home
    \item \textbf{Screen use:} digital screens cause eye strain and discomfort
    \item \textbf{Heredity:} many eye conditions run in families
\end{itemize}

\section*{Clinical Features}
Good eye care involves knowing the warning signs.
\begin{itemize}
    \item Blurred, distorted, or reduced vision
    \item Eye pain, redness, or discharge
    \item Floaters, flashes, or a shadow in the vision
    \item Difficulty seeing at night
    \item Frequent squinting or headaches
    \item In children: poor school performance or an eye turning
    \item Changes in vision after the age of 40
\end{itemize}

\section*{Differential Diagnosis}
Eye symptoms have many causes.
\begin{itemize}
    \item Dry eye, infections, and allergies cause redness and irritation
    \item Migraine and other conditions cause visual disturbance
    \item Stroke and other neurological conditions can affect vision
    \item Serious conditions such as retinal detachment and glaucoma require urgent care
    \item Prompt assessment distinguishes minor from serious causes
\end{itemize}

\section*{When to Seek Medical Care}
Everyone should have regular eye examinations: children at school age, adults every two years, and more often with age, diabetes, or eye disease. Any new or sudden vision symptom requires assessment.

\textbf{Contact your local emergency services for:}
\begin{itemize}
    \item Sudden loss of vision
    \item Sudden severe eye pain
    \item A red eye with severe pain and nausea
    \item Sudden floaters with flashes, which may indicate retinal detachment
    \item A chemical in the eye or a penetrating eye injury
\end{itemize}

\section*{Diagnosis and Investigations}
Eye examinations detect problems early.
\begin{itemize}
    \item \textbf{Visual acuity:} the eye chart test
    \item \textbf{Refraction:} determines the prescription for glasses
    \item \textbf{Slit-lamp examination:} inspects the front of the eye
    \item \textbf{Eye pressure:} checks for glaucoma
    \item \textbf{Dilated examination:} views the retina and optic nerve
    \item \textbf{Screening for diabetes:} regular retinal photography for people with diabetes
\end{itemize}

\section*{Management}
Eye care is managed with prevention, correction, and treatment.
\begin{itemize}
    \item \textbf{Glasses and contact lenses:} correct refractive error
    \item \textbf{Protection:} sunglasses, safety glasses, and sports goggles
    \item \textbf{Eye hygiene:} good practices for contacts and around screens
    \item \textbf{Treatment of conditions:} drops, laser, and surgery for cataract, glaucoma, and other conditions
    \item \textbf{Managing diabetes:} good blood sugar control protects the eyes
    \item \textbf{Regular reviews:} monitoring for people at risk
\end{itemize}

\section*{Complications}
\begin{itemize}
    \item Avoidable vision loss from undetected or untreated conditions
    \item Blindness from glaucoma, diabetes, and retinal disease
    \item Eye injuries from unsafe work, sport, or handling chemicals
    \item Falls and injuries from poor vision in older people
    \item Reduced quality of life and independence
\end{itemize}

\section*{Prognosis and Outcomes}
Most vision loss is preventable or treatable when detected early. Regular examinations catch glaucoma, macular degeneration, diabetic eye disease, and cataract at treatable stages, and eye protection prevents many injuries. Good eye care throughout life keeps vision clear and preserves independence.

\section*{Prevention and Self-Care}
\begin{itemize}
    \item Have regular eye examinations at every age
    \item Wear sunglasses that block UV light
    \item Use protective eyewear for work, sport, and DIY
    \item Follow the 20-20-20 rule for screens
    \item Manage diabetes, blood pressure, and weight
    \item Do not smoke, and eat a diet with plenty of fruit and vegetables
    \item Seek prompt care for any eye symptom
\end{itemize}

\section*{Sources and Further Reading}
The following reputable sources provide further information for healthcare students and patients seeking reliable, current advice about eye care.
\begin{itemize}
    \item Optometry Australia. \textit{Eye care and vision health.} https://www.optometry.org.au/
    \item The Royal Australian and New Zealand College of Ophthalmologists. \textit{Eye health information.} https://ranzco.edu/
    \item Vision Initiative. \textit{Preventing avoidable blindness.} https://www.visioninitiative.org.au/
    \item Healthdirect Australia. \textit{Eye health.} https://www.healthdirect.gov.au/eye-health
    \item NHS. \textit{Eye health and vision.} https://www.nhs.uk/conditions/
    \item MSD Manual. \textit{Eye disorders.} https://www.msdmanuals.com/
\end{itemize}
